\chapter{70}



Louis trat in die Hofküche, blieb stehen und sah Mazi beim Einräumen der eben
gelieferten Einkäufe zu. Mazi blickte ihn gutgelaunt an.

«Kaffee?»

«Gute Idee, gerne.»

An die Wand gelehnt steckte sich Louis die Hände in die Hosentaschen und schaute Mazi beim
Kaffeemahlen zu. Der musste gleich
mehrere Glücksgene geerbt haben. Mit Schwung stellte Mazi zwei Espressi
auf den Tisch. Er schob die eine Tasse vor Louis hin und trank seine mit zwei
Schlucken leer.

Unterdessen kam sich Louis vor, als sässe er auf einer Schaukel. Er
pendelte hin- und her, um im geforderten Moment abzuspringen. Von einer
Realität in die andere. Und er wunderte sich, wie die Wechsel
unterdessen von sich gingen. Wie ihm noch immer gelang, die beiden
Lebensbereiche voneinander zu trennen. Louis und Ciro waren
höchstens entfernte Verwandte, die in unterschiedlichen Welten lebten
und sich gegenseitig ausschliessen mussten. Wie lange es wohl noch gutgehen konnte,
eine Vermischung zu verhindern? 

«Wie geht’s dir?»

Louis fuhr zusammen. Mazi setzte sich und verschränkte die Arme. 

«Gut, gut, ich habe nur gerade an mein neues Stück gedacht. Ich komme
nicht weiter.»

Er unterstrich seine Worte mit einem Stirnrunzeln.

«Verstehe, also, wenn du Zeit für dich brauchst, übernehme ich die
Reinigung der Praxis alleine, kein Problem, draussen bin ich sozusagen
fertig, hm?»

«Bist du sicher?»

«Aber natürlich, Ciro, wenn es dir hilft.»

«Danke, Mazi, grosszügig von dir, vielleicht am Nachmittag eine Stunde?»

Mazi stand auf.

«Klar,» er räumte den Tisch ab, «kochen wir heute zusammen? Ich glaube,
ich kann da so einiges von dir lernen.»

Während Mazi ihm zuzwinkerte, fiel Louis ein, wie der ihm von seiner Herkunft erzählt hatte.
Von seinen Eltern, deren Verhältnis zueinander und zu ihm als Sohn.
Louis’ Geschichte erzählte sich ungleich anders. Lucia und er waren
Kinder aus der Kombination zweier Menschen, die unterschiedlicher nicht
sein konnten. Eine Fehlkomposition? Der Gedanke war merkwürdig. Wie
weit zurück ihn seine Erinnerungen nur immer wieder trugen. Es war
schräg.

\sectionbreak

\indentedblock{Damals. Fribourg, acht.\\
Louis trommelt mit den Fingerspitzen auf dem Fensterbrett herum. Er
wartet. Lucia hört er aus dem Nachbarzimmer mit ihren Plüschtieren
sprechen. Schon lange. Und wie immer. Das Geplapper geht ihm auf die
Nerven. Ihr scheint es egal zu sein, dass ihr Vater immer noch nicht da
ist. Jetzt warten sie schon über eine Stunde. Maman meint, er kommt
bestimmt. Vielleicht wird er aufgehalten. Die Uhr dreht weiter, Minute
um Minute. Mamans Handy klingelt. Louis hört sie reden. Leise, gedämpft
nur, wie oft, wenn sie nicht will, dass er mithört. Er weiss es schon
lange. Sie streiten wieder. Und es geht um ihn. Manchmal auch um Lucia.
Aber mehrheitlich um ihn.

Kurz ist es still. Dann hört er Mamans Schritte. Sie kommt in sein
Zimmer und setzt sich auf sein Bett. Sie macht ein bekümmertes Gesicht.
Louis weiss, was kommt. Sie könnte sich die Worte sparen. Louis spürt,
wie die Wut kommt.

«Papa hat noch zu tun, ein neuer Kunde, Louis. Nächsten Samstag holt er
euch bestimmt ab. Komm her.»

Mamans Stimme hat den bekannten Ton. Sie breitet die Arme aus. Sie weiss, dass
er enttäuscht ist. Sie will ihn trösten.

«Mir reichts,» schreit Louis, «warum muss ich euer Kind sein. Warum habt
ihr eigentlich Lucia und mich gemacht, warum?»

Louis wirft sich aufs Bett und zieht die Decke über den Kopf. Jetzt
weint er, obwohl er nur wütend sein will. Dann schreit er die Worte
durch die Decke ins Kinderzimmer.

«Du hast Papa rausgeworfen, Maman, du bist Schuld. Du hast unser Leben
kaputt gemacht. Papa hätte das nie getan. Er hat mir alles erzählt. Ich
will zu Papa.»

Maman sitzt ruhig da und hört Louis zu. Zwischendurch dringt Lucias
Elefantenstimme über den Flur zu ihnen hinüber. Louis kann sie nicht
mehr hören. Er reisst die Decke hinunter.

«Sei endlich still da drüben,» brüllt er.

Maman berührt nun Louis’ Schulter. Er schluchzt und lässt es geschehen.
Sie nimmt ihn in die Arme. Louis sinkt in sich zusammen. Wie immer.
Jetzt steht Lucia im Türrahmen. Maman winkt sie zu sich. Sie kuschelt
sich dazu.

«Wir haben uns schnell verliebt, euer Papa und ich. Ihr wisst, wir haben
uns in Spanien am Meer das erste Mal gesehen. Er hat mir so gut
gefallen.»

Louis schaut hoch. Er liebt die ganz alten Maman-Papa-Geschichten und
wischt sich die Tränen aus den Augen. Maman sieht abwesend aus. Sie
schaut aus dem Fenster.

«Er war anders als die andern Männer. Er kam mit einem gelben Kabrio
angebraust. Ich sass im Strandcafé mit Louise…»

«Louise?»

«Ja, sie war schon damals meine Freundin.»

Maman lächelt. Lucia drückt sich tiefer in Mamans Arm.

«Wir hatten grosse Pläne. Und dann kamst du zu uns, Louis. Du konntest
es wohl kaum erwarten,» lacht sie und Louis schmunzelt.

«Und dann du, Lucia,» Maman kitzelt die Kleine, «ihr seid das
Allerbeste, was uns passieren konnte. Ich liebe euch.»

Jetzt sieht Louis ganz kleine Tröpfchen in Mamans Augen. Er mag nicht,
wenn sie weint.

«Ja, und mit den Jahren wurde es schwierig. Uns haben nicht mehr die
gleichen Dinge gefallen. Papa wollte nach Spanien in die Ferien,» sie
macht eine Handbewegung nach links, «ich lieber nach Frankreich,»
Handbewegung nach rechts, oder - » Maman denkt nach, «oder wir konnten
uns nicht einigen, wofür wir unser Geld ausgeben wollten. Papa liebt
teure, schnelle Autos und ich wünschte mir ein kleines Haus, eins, wie
es Grandpère hat. Wir haben immer mehr gestritten und das war nicht gut,
auch für euch zwei nicht.»

Louis setzt sich auf. Er sieht Madame Aebischer vor sich. Seine
Lehrerin. Es fühlt sich gut an. Er erzählt, wie Streit in seiner Schule
geschlichtet wird.

«…und dann könnt ihr Frieden machen, Maman, euch die Hand geben und
gut ist.»

Aber noch während er redet, weiss er, dass das nichts wird. Maman will
das nicht.

«Wir haben vieles versucht, Louis. Es hat nicht geklappt. Aber eure
Eltern bleiben wir ein ganzes Leben lang, versprochen.»

Maman hat jetzt diesen lieben Blick. Dann hat sie eine ihrer guten
Ideen. Sie packen in Windeseile ihre Badesachen und fahren ins
Schwimmbad.

«Nein, leider nicht,» hört Louis Maman leise sagen, bevor sie ins Auto
einsteigt, sie ist am Handy, «er hat abgesagt - ja, ich weiss, es tut
mir leid - ja, ich danke dir, vielleicht nächstes Wochenende. Danke, das
ist sehr lieb von dir. Ciao.»

«Wer war’s?» fragt Louis.

«Marie. Alles in Ordnung, Louis.»

An diesem Tag wagt Louis seinen ersten Sprung vom kleinen Turm. Lucia
und Maman stehen am Beckenrand. Sie klatschen in die Hände, als er
auftaucht. Er fühlt sich wie ein Mann.

Er wird seinem Vater erzählen, wie mutig er war.}
